\subsubsection{The ordinary CM chart and its theta measure}
\label{subsec:word-measure}

The moments depend on more than the moduli of the theta coefficients. Translation on the ordinary CM chart records the character and support of the measure used in both tests.

For the ordinary chart use the imaginary quadratic order
\[
 K_0=\Q(\sqrt{-d}),\qquad O=\Z[(1+\sqrt{-d})/2],
 \qquad d\text{ prime},\quad d\equiv7\pmod{16}.
\] In this section \(p\)
is the odd split prime assigned to the sequence and \(\lambda\) is the chosen embedding
of the algebraic coefficients in \(\mathbb C_p\). We use subscripts
\(+,-\) for the two \(p\)-adic components of \(O\), with the first
induced by \(\lambda\). We may transport formal complex Taylor
coefficients by a field isomorphism \(\iota:\C\simeq\mathbb C_p\)
extending \(\lambda\) on \(\overline{\Q}\). Both fields are
algebraically closed of transcendence degree continuum over
\(\overline{\Q}\), so a bijection of transcendence bases extends
the prescribed action on the algebraic coefficients. This transport
uses the field operations on Taylor coefficients.

For the order lattice \(O\) put
\(H(z,w)=2z\bar w/\sqrt d\), and write
\[
\begin{split}
 U_wF(z)&=e^{-\pi H(z,w)-\pi H(w,w)/2}F(z+w),\\
 W_wF&=e^{\pi i |w|^2}U_wF,\qquad
 \kappa(v,w)=e^{-2\pi i\operatorname{Im}H(v,w)}.
\end{split}
\]
Let \(\theta\) be the degree-one even theta vector with
\(\theta(0)=1\), \(W_l\theta=\theta\) for \(l\in O\), and zero
at \(h=1/2\) modulo \(O\). These data follow from the
degree-one theta transformation law. In the basis
\(1,(1+\sqrt{-d})/2\) the quadratic sign
\(e^{\pi i|l|^2}\) is \((-1)^{a+ab}\) on \(l=a+b(1+\sqrt{-d})/2\).
Here the norm of the second generator is \((1+d)/4\), which is
even. The theta characteristic is the translate of the odd
theta section by \(1/2\). The divisor has degree one, so
this translation gives its only zero.

\begin{proposition}[The ordinary chart with its chosen torsion coordinate]\label{prop:word-ordinary-chart}
At the assigned split prime, the ordinary formal coordinate
has transported logarithm \(\omega\log(1+X)\) with
\(|\omega\iota(\delta_0)|_p=1\), where \(\delta_0^{12}=\Delta(O)\).
Choose it so that at a connected torsion point \(y=v/p^n\)
with \(v_-\in p^{2n}\Z_p\), the coordinate specializes according to
\[
 1+X(y)=\iota(\kappa(1,v/p^n)).
\]

\end{proposition}

\begin{proof}
\proofstep{Obtain a good ordinary formal model and its unit differential.}

Write \(\delta_0^{12}=\Delta(O)\) for the discriminant in
the analytic uniformization by \(\C/O\). The CM \(j\)-invariant
is an algebraic integer, so after base extension the elliptic
curve has good reduction. A good Weierstrass model can be
chosen over algebraic coefficients. With analytic scaling
\(c\) from the lattice invariants to this model,
\(\abs{\iota(\delta_0)/\iota(c)}_p=1\), by the discriminant
transformation of weight twelve. Good reduction is ordinary
here. Indeed the two idempotents of

\[
  O\otimes\Z_p=\Z_p\oplus\Z_p
\]
 give the two height-one
parts of the \(p\)-divisible group and the tangent map chooses
the first. Over the completed algebraic closure the connected
part is \(\widehat{\mathbb G}_m\). Its isomorphism with the
formal elliptic curve is integral and has unit linear
coefficient in a good parameter at the identity. It can be
obtained by choosing a compatible generator of the dual
\'{e}tale \(p\)-power group and using Cartier duality at every
level. Write \(z(X)\) for the complex logarithmic coordinate after this transport. On the multiplicative formal group the resulting logarithm is
\[
  z(X)=\omega\log(1+X),\qquad z'(0)=\omega.
\]
 The two
invariant differentials differ by a unit in good coordinates,
so \(\Omega_0=\omega\iota(\delta_0)\) has norm one.
\cite{bannai-kobayashi-theta}\label{cor:word-measure-1}

\proofstep{Normalize by the connected torsion pairing.}

For \(y=v/p^n\) with \(v\in O\) and
\(v_-\in p^{2n}\Z_p\), \(y\) belongs to the connected
\(p^n\)-torsion. Its coordinate \(1+X\) is
\(\zeta_y=\iota(\kappa(1,y))\). Indeed on that line the pairing
with \(1\) is a generator of its dual at every level.

In the basis \(1,(1+\sqrt{-d})/2\), if \(v=a+b(1+\sqrt{-d})/2\),
then 
\[
  \kappa(1,v/p^n)=e^{2\pi i b/p^n},
\]
 and \(b\) modulo
\(p^n\) is a unit multiple of \(v_+\) on the connected factor.
The ratio between this identification and any multiplicative
coordinate on the connected torsion is multiplication by a
compatible system of units modulo \(p^n\). Replacing the
initial coordinate by its power for the resulting unit of
\(\Z_p\) enforces this normalization without changing
\(\abs{\Omega_0}_p\). All connected torsion has lifts with the
stated extra divisibility by the Chinese remainder theorem.

\end{proof}

For a complex germ \(G\), let \(\widehat G\) be the series
obtained by applying \(\iota\) to its Taylor coefficients
and substituting \(\omega\log(1+X)\). The coefficient
supremum norm of a bounded series is its supremum on the open
unit disc. This follows by applying the Gauss norm at radii
tending to one. The norm is multiplicative and is unchanged
by the substitutions
\[
 h_a(X)=(1+X)^a-1\quad(a\in\Z_p^\times),\qquad
 X\longmapsto \zeta(1+X)-1\quad(\abs{\zeta-1}_p<1).
\]
Each substitution maps the open unit disc into itself
in both directions. We also use the following fact about an
elliptic rational function with no pole in an identity disc
or another specified residue disc of a good model. The two
expansions in integral parameters have the same finite
coefficient norm.
Indeed, if the two parameters are related by \(\phi\) and \(\phi^{-1}\), both integral, the coefficient norm comparison in both directions gives
\[
  \|F\circ\phi\|\le\|F\|=\|(F\circ\phi)\circ\phi^{-1}\|\le\|F\circ\phi\|.
\]
 This is the norm at the generic point of the
special fiber.

To see the assertion at the identity, multiply
the rational function by a monic polynomial in the Weierstrass
coordinate \(x\) whose roots are the \(x\)-coordinates of all
its poles off the identity, with enough multiplicity to clear
those poles. Each root is integral when the pole is outside the
identity disc. The resulting numerator is a polynomial in
\(x,y\), linear in \(y\). Its monomials have distinct pole
orders at the identity.
Indeed in a Weierstrass parameter at the identity \(O\) the two lists of orders are
\[
  \operatorname{ord}_O(x^i)=-2i,\qquad \operatorname{ord}_O(x^iy)=-2i-3\qquad(i\ge0).
\]
 Multiplication by the corresponding
power of the formal parameter makes each expansion integral
with leading coefficient a unit. After scaling by a coefficient
of maximum norm, the terms of maximum norm have a nonzero
reduction, as their initial orders are distinct. The monic
denominator, shifted by its pole order, is a power series
unit.

For its shifted denominator \(D(T)\) this uses the unit criterion
\[
  D(0)\in{\cal O}^{\times}\quad\Longrightarrow\quad D(T)^{-1}\in{\cal O}[\![T]\!].
\]
 This computes the norm of the expansion. Equivalently,
the numerator after this scaling has nonzero reduction as a
rational function on the special fiber and the denominator is
monic in \(x\). Translating by a point of the good model
preserves that generic norm and proves the assertion at the
other disc. The expansion of the translated rational function
at a point of the identity disc is therefore its analytic
translation on that disc. These arguments apply to the
transported elliptic quotients by expressing them as rational
functions of the Weierstrass coordinates on the complex curve
and transporting the finite rational expression.

\paragraph{Integral translations of the degree-one theta vector.}

The ordinary local measure is built from translated theta vectors. Begin with translations of order prime to the residue characteristic, where the coefficient norm is preserved as asserted below.

\begin{proposition}[Coefficient norms of prime-to-\(p\) translations]\label{prop:word-integral-translates}
In the ordinary chart, the normalized theta vector and
each prime-to-\(p\) translate have coefficient norm at
most one. For a finite sum \(G\) of fixed powers of these translates and any allowed translation by \(t\), the coefficient norm satisfies
\[
 \|U_tG\|_{\rm coeff}=\|G\|_{\rm coeff}.
\]
 This
also holds coordinatewise for sums of products.
\end{proposition}

\begin{proof}
\proofstep{Use a CM norm divisible by \(p\).}

Choose \(\alpha\in O\) with \(\alpha_+\) a unit and
\(\alpha_-\equiv0\pmod p\), and set \(m=\alpha\bar\alpha\).
The quotient
\[
 R_\alpha(z)=\theta(\alpha z)/\theta(z)^m
\]
is elliptic, since the two factors have identical automorphy
on \(O\). It and its inverse have no poles in the identity
disc. A zero of either theta factor would map under \(1\)
or \(\alpha\) to the nonzero two-torsion \(h\). Its
reduction is nonzero because \(p\) is odd, whereas an
endomorphism preserves the identity disc.

Since \(R_\alpha(0)=1\),
its series and inverse are bounded, and multiplicativity gives
norm one for both. Consequently the formal equation for
\(F=\widehat\theta\) is
\[
                  F(h_{\alpha_+}(X))=F(X)^m\widehat R_\alpha(X).
\tag{M1}\label{eq:tag-m1}
\]
The coefficient of degree \(j>0\) determines the coefficient
of \(F\) there with unit multiplier \(\alpha_+^j-m\).
The terms of smaller degree are integral by induction, from
\(F(0)=1\). This proves \(\|F\|=\|F^{-1}\|=1\).

\proofstep{Apply the isogeny on a prime-to-\(p\) residue disc.}

Suppose \(t\in K_0\) satisfies \(s t\in O\) for some positive
\(s\) prime to \(p\), and \(t\not\equiv h\pmod O\).
Choose \(\alpha\in O\) with \(\alpha_+\) a unit,
\(\alpha_-\equiv0\pmod p\) and \(\alpha\equiv1\bmod sO\).
Using \(H(\alpha z,\alpha w)=mH(z,w)\) gives
\[
 U_{\alpha t}\theta(\alpha z)=(U_t\theta(z))^m R_\alpha(z+t).
\]

Since \(\alpha t-t\in O\), \(U_{\alpha t}\theta\) differs from
\(U_t\theta\) by a root of unity. Both \(R_\alpha\) and its
inverse are regular throughout the disc of \(t\), since the
reductions of both \(t\) and \(\alpha t=t\) on the curve
differ from that of \(h\). Their translated series have
norm one by the generic norm comparison. In particular the
constant term of the translated quotient is a unit. The
nonzero constant term of \(\widehat{U_t\theta}\) has norm one
by the constant term of the displayed functional equation.
Dividing by this constant gives \eqref{eq:tag-m1} with an integral
quotient of constant term one. The same induction gives a
unit series for \(\widehat{U_t\theta}\).

\proofstep{Compare the norm in the two residue discs.}

Fix a power \(a\) and write the finite sum as
\[
  G=\sum_j c_j(U_{t_j}\theta)^a.
\]
Choose a common integer \(s\), prime to \(p\), clearing the denominators of \(t\) and of all \(t_j\).
The function
\[
               R(z)=G(sz)/\theta(sz)^a
\]
is elliptic. Its potential poles are where \(sz\equiv h\).
If \(t\not\equiv h\), the discs at \(0\) and at \(t/s\) avoid
these poles, since \([s]\) is prime to \(p\) and \(t\) and
\(h\) are separated in reduction. Its coefficient norm is
\(\|\widehat G\|\) there, as \(\widehat\theta\) and its inverse
have norm one. From
\[
       \sum_j c_j (U_t(U_{t_j}\theta)(sz))^a
                =(U_t\theta(sz))^a R(z+t/s)
\]
and the preceding unit-series bound, translation by \(t\)
preserves this norm.

When \(t\equiv h\), apply the bound
twice for \(t/2\), which is prime to \(p\) and is not the
half-period modulo \(O\). The two translations along
\(t/2\) compose to \(U_t\). In particular every prime-to-\(p\)
translate of \(\theta\) has norm at most one. Applying the
same proof one variable at a time gives the norm preservation
for sums of products of fixed powers in several variables.
\label{cor:word-measure-2}

\end{proof}

\paragraph{Connected torsion acts through the commutator.}
\begin{proposition}[Connected torsion on the same ordinary series]\label{prop:word-connected-torsion}
Take a connected torsion lift \(y=v/p^n\) with
\(v\in O\) and \(v_-\in p^{2n}\mathbb Z_p\), and a translate \(t\)
of order prime to \(p\). In the coordinates of Proposition~\ref{prop:word-ordinary-chart}, the transformation law \textup{\eqref{eq:tag-m2}} is
\[
 \widehat{W_y(U_t\theta)}(X)=\widehat{U_t\theta}\big(\zeta_y(1+X)-1\big).
\]
 Its derivatives
in the stipulated \(\delta_0^{-1}\)-scaled directions are
algebraic and retain the unit bound.
\end{proposition}

\begin{proof}
\proofstep{Use the commutators to determine the scalar.}

Write \(y=v/p^n\) for this lift. For \(\alpha_+\equiv1\pmod{p^n}\)
we have \(\alpha y-y\in O\), and \(y\bar y\in\Z\) by the
condition on \(v_-\). The quadratic factor and the
Heisenberg composition of the \(W\)'s then give
\(W_{\alpha y}\theta=W_y\theta\): the product
\(y\overline{(\alpha-1)y}\) is in \(O\), on which the
quadratic cocycle is integral. Put
\(P(X)=F(\zeta_y(1+X)-1)\) and \(Q=\widehat{W_y\theta}\).
The dilation \(h_{\alpha_+}\) commutes with the indicated
substitution because \(\zeta_y^{\alpha_+}=\zeta_y\).

Translating \eqref{eq:tag-m1} and transporting the identity for
\(W_{\alpha y}\theta(\alpha z)\) gives
\[
 (Q/P)(h_{\alpha_+}(X))=(Q/P)(X)^m .
\]
Here the translated quotient is \(R_\alpha(z+y)\) on
both sides, and \(P\) is a unit series. The ratio is nonzero
since \(W_y\theta\) is nonzero. Comparing orders at \(X=0\)
shows that the ratio has order zero. Its constant term
\(b\) satisfies \(b^{m-1}=1\), and comparing the first possible
nonconstant coefficient with multiplier \(\alpha_+^j-m\)
shows the ratio is constant. This holds for all such
\(\alpha\) with \(p\mid m\). First take one of norm \(m_1\).
Then additionally require \(m_1-1\mid\alpha\) in \(O\),
which is compatible with the \(p\)-conditions. The resulting
norm \(m_2\) is divisible by \(m_1-1\), so

\[
  \gcd(m_1-1,m_2-1)=1
\]
 and \(b=1\).

\proofstep{Translate the elliptic quotient in the connected chart.}

For \(U_t\theta\), take \(s\) prime to \(p\) with \(st\in O\)
and a primary lift \(y=s v/p^n\) with
\(v_-\in p^{2n}\Z_p\). Use the elliptic quotient
\(R(z)=U_t\theta(sz)/\theta(sz)\).
On replacing \(z\) by \(z+y/s\) the chart expression
of \(R\) transforms by
\(1+X\mapsto\zeta_y^{1/s}(1+X)\), since \(y/s\)
is the indicated lift of connected torsion.
Here \(1/s\) in the root of unity is taken modulo \(p^n\).

Since \(W_y(U_t\theta)(sz)=W_y\theta(sz)R(z+y/s)\),
cancellation of the invertible substitution \(h_s\) gives
\[
    \widehat{W_y(U_t\theta)}(X)
              =\widehat{U_t\theta}(\zeta_y(1+X)-1).
\tag{M2}\label{eq:tag-m2}
\]
The elliptic quotient has no pole in the identity disc,
so the rational comparison applies there. The proof also
applies for \(t\equiv h\), since the quotient uses
\(\theta(sz)\) in the denominator.

\proofstep{Check the coefficient transport for the scaled derivatives.}

The jets in \eqref{eq:tag-m2} in directions scaled by \(\delta_0^{-1}\)
are algebraic. Indeed \((1+X)d/dX\) preserves the unit
coefficient bound. Applied repeatedly at zero it gives the
transported derivatives, with powers of \(\Omega_0\) of norm
one. The same bound holds for every field isomorphism
extending \(\lambda\), with its own multiplicative chart and
unit period. If one of the scaled derivatives were
transcendental, an automorphism of \(\C\) fixing
\(\overline{\Q}\) could multiply it by \(p^{-k}\), by
extending from a transcendence basis. The absolute values
of its images would then be unbounded. This proves the
algebraicity needed for transport by \(\lambda\).
\label{cor:word-measure-3}

\end{proof}

\paragraph{Measures.}

The coefficients of the ordinary power series are the binomial moments of its measure. Translation and finite character projection can therefore be checked on these moments.

A bounded power series \(\sum_\nu a_\nu X^\nu\) in \(g\)
variables determines a measure \(\mu\) on \(\Z_p^g\) by

\[
  \mu(\prod_i\binom{A_i}{\nu_i})=a_\nu,
\]
 with norm
\(\sup_\nu\abs{a_\nu}_p\). The product Mahler basis is
orthonormal for continuous functions on \(\Z_p^g\), so the
measure evaluates a Mahler expansion by a convergent sum.
The generating series is \(\mu_A(\prod_i(1+X_i)^{A_i})\).
In this representation \((1+X_i)d/dX_i\) multiplies by
\(A_i\), and substitution

\[
  1+X_i\mapsto\zeta_i(1+X_i)
\]
 multiplies by
\(\zeta_i^{A_i\bmod p^n}\) when \(\zeta_i^{p^n}=1\).
Indeed these identities hold for each integer sample and
extend to \(\Z_p^g\) by continuity, with the substitutions
converging on every smaller disc. Thus \eqref{eq:tag-m2} realizes
connected torsion translates as twists of the same measure.
\cite{bannai-kobayashi-theta}\label{cor:word-measure-4}

\paragraph{The Hom period lattice.}
In a Hom space of the selected kernel let \(M\) be its period
lattice in complex coordinates, and let \(H\) be its
Hermitian form. Use
\(\beta=\sqrt d\,H/2\), so that
\(\operatorname{Im}H\) is the alternating form that is
perfect on \(M\) and \(\beta(M,M)\subset O\).
The same formulas for \(U,W,\kappa\) apply, with
\(\beta(y,y)\) in place of \(|y|^2\).
\begin{proposition}[The ordinary bound in the actual Hom lattice]\label{prop:word-hom-lattice-integral}
For \(e^\natural\) of the selected actual Hom pair and
any prime-to-\(p\) translate, write \(\sum_\alpha b_\alpha\mathbf X^\alpha\) for the series in a finite list of actual period directions in the ordinary chart. Then
\[
 \sup_\alpha|b_\alpha|_p\le1.
\]
 The logarithmic scale \(\omega\)
is common to these directions and to the standard pair.
\end{proposition}

\begin{proof}
\proofstep{Test along the actual period lattice.}

For \(v_1,\ldots,v_k\in M\) test the series
\[
      G\!\left(\sum_{j=1}^k
             \omega\log(1+X_j)v_j\right)
\]
after the specified field transport. Bound one means that
all its coefficients belong to the valuation ring. Such
series are closed under products and integral linear
combinations. If a linear change sends a direction \(v_j\)
into the lattice after multiplying it by an integer \(s\)
prime to \(p\), use the series for \(s v_j\) after substituting
\(X_j\mapsto h_{1/s}(X_j)=(1+X_j)^{1/s}-1\).
The binomial coefficients are integral since \(1/s\in\mathbb Z_p\),
and the constant term is zero. The identity

\[
  \log(1+h_{1/s}(X_j))=s^{-1}\log(1+X_j)
\]
 gives the claimed
pullback and preserves bound one.

\proofstep{Compare with the good model by the prime-to-\(p\) isogeny.}

For the base curve lattice \(\Gamma=\delta\mathfrak c\),
the ideal \(\mathfrak c\) has norm prime to \(2p\).
The isogeny \(\C/\mathfrak c\to\C/O\) has order prime to
\(p\) and identifies the two formal groups on good models.
In good coordinates the two invariant differentials have
unit ratio under this isogeny. In coordinates for the
lattices without the scale \(\delta\) their differentials
are matched by the identity map. In particular
\(\delta/\delta_0\) has \(p\)-adic norm one.

In valuations this comparison is
\[
  v_p(\delta/\delta_0)=v_p(\delta)-v_p(\delta_0)=0 .
\]
 This follows
also by the equality of the valuations of the two
discriminants in their analytic coordinates, since each
normalized discriminant is a unit and the isogeny has
unit differential. The equality holds for every field
transport extending \(\lambda\), which proves
\(\delta/\delta_0\in\overline{\Q}\) by the same
automorphism argument used for the theta jets.
The half-period of the even theta section on \(\Gamma\)
is nonzero of order two by the parity computation in
Section~\ref{subsec:theta-normalization}.
Proposition~\ref{prop:word-integral-translates} applies to
that section on \(\Gamma\): the endomorphisms are still
\(O\), the same automorphy yields the elliptic ratios,
and the half-period is still separated from zero in
reduction. The pulled-back multiplicative chart now uses
logarithmic scale \(\iota(\delta)\omega\). For the
lattice-direction series above, each direction

\[
  v_j=\delta b_j
\]
 uses the substitution associated to
the integral exponent \(\lambda(b_j)\in\Z_p\).
Writing \(b_j\) for these directions in the unscaled
coordinates, the multiplicative parameter of their sum is
\[
           \prod_j(1+X_j)^{\lambda(b_j)}-1 .
\]
It has integral coefficients and zero constant. Composing
the integral one-variable section with this parameter proves
bound one for the prime-to-\(p\) translates in all the
chosen directions.

\proofstep{Pull back the orthogonal standard chart.}

Apply the fourth-tensor normalization of
Section~\ref{subsec:theta-normalization} to the
base CM vector. At \(p\) the rational isometry of the four
copies sends the period lattice to the standard curve
lattice. Its global intersection indices are prime to
\(p\). The vector after an arbitrary prime-to-\(p\)
translation becomes a finite sum of products of
one-variable translates. Every coefficient in that sum
is a product of the normalized unit, an inverse square
root of a prime-to-\(p\) index and Heisenberg roots of unity.
For an index \(I\) prime to \(p\) and a root of unity \(\zeta\), the two frame factors contribute
\[
  v_p(I^{-1/2}\zeta)=-\tfrac12v_p(I)+v_p(\zeta)=0 .
\]
Every translation in its product factors has order prime
to \(p\). Pull back the bound one from the standard
directions, multiplying by an intersection index
before the pullback if necessary.

Write \(A(X)=\sum_\alpha a_\alpha X^\alpha\) for the one-copy
reading after this pullback. The four test blocks are independent.
Independence of the four variable blocks gives
\[
  [X_1^\alpha X_2^\alpha X_3^\alpha X_4^\alpha]\prod_{\nu=1}^4 A(X_\nu)=a_\alpha^4.
\]
The bound for the fourth tensor gives \(4v_p(a_\alpha)\ge0\)
for every nonzero coefficient, without needing a nonzero
constant term. Tensoring these one-variable estimates along
the matrix entries gives the same bound for the standard Hom
vector and the same scale \(\omega\).

\proofstep{Use the two period-lattice averages in the actual frames.}

For general source and target in the Hom space,
rational unitary changes put both lattices in standard
position at \(p\). This is the split lattice argument
in Section~\ref{subsec:theta-normalization}, using a basis
and its dual at \(p\) and weak approximation on the Cayley
chart. The basic theta identity for the selected vector
uses the two normalized period-lattice averages,
first changing the source, then the target. After
these changes the intermediate lattices and the standard
ones agree at \(p\). For either average, write \(\mathcal D=M_2/(M_1\cap M_2)\).
The normalized translation of the section is
\[
      |\mathcal D|^{-1/2}\sum_{a\in\mathcal D}W(a)G .
\]

In our two steps its normalization is a local unit, because
the intermediate and standard lattices agree at \(p\).
This is the theta-line transfer of Section~\ref{subsec:theta-normalization}.
Here both such indices are prime to \(p\).
The determinant character between the transported
combinations has norm one, as its determinant ideals
agree at \(p\) and the remaining finite characters
have finite image. The resulting coefficients are
integral. Adding a prime-to-\(p\) translate preserves
the order condition on every summand by the lattice
intersection inclusion. Thus the bound one holds for
the actual normalized vector \(e^\natural\) and all
its prime-to-\(p\) translates. The rational changes
carry both the vector and its argument lattice, so
the test directions are the actual period directions.
\label{cor:word-measure-5}

\end{proof}

\begin{proposition}[An integral theta measure for decorated jets]\label{prop:theta-measure-word}
Use the ordinary model at the assigned split prime and
the normalized vector on the actual Hom lattice above.
Write \(\mu\) for the measure in the selected coordinates. For a monomial moment test \(P\), its theta jet value \(J(P)\) satisfies
\[
 J(P)=\int P\,d\mu,\qquad \|\mu\|\le1.
\]

A finite list of those moments is approximated at any
prescribed precision by one set of samples in the
rational Hom period lattice with integral coefficients.
\end{proposition}

\begin{proof}
Choose a \(K_0\)-basis \(x_j\) of the rational space of \(M\)
with \(x_j\in M\), orthogonal for \(\beta\), and with
\(\beta(x_j,x_j)=a_j\in\Z_{>0}\) prime to \(p\).
To find one, take a basis and the dual basis in the two
split factors of the unimodular lattice, approximate
both modulo \(p\) by rational vectors of \(M\), and
orthogonalize. The pivots are units at \(p\), so the
denominators can be cleared by integers prime to \(p\).
The resulting \(x_j\) remain a basis after localization
at \(p\). The orthogonal frame determines the pair of lattices
\[
 M_0=\bigoplus_j O x_j\ \subseteq\ M\ \subseteq\ M_0^\#
        =\bigoplus_j a_j^{-1}O x_j ,
\]
where \(\#\) denotes the alternating dual. The section on the smaller
lattice factors in this frame as \(F_0(z)=\prod_j\theta(z_j)^{a_j}\).

For any \(t_0\in M_0^\#\),
\(U_{t_0}F_0\) is fixed by the Weyl operators of \(M_0\).
Summing its Weyl translates over \(M/M_0\) gives a theta
vector for \(M\). At least one such sum is nonzero:
evaluate at \(-t_0\) and divide by the common
nonzero Gaussian there. The coefficients of the
characters 
\[
  t_0\mapsto\kappa(l,t_0)
\]
 for \(l\in M/M_0\)
are then the corresponding Weyl translates of \(F_0\) at
zero, with coefficient one for \(l=0\).
These characters are distinct by duality. The theta
space has dimension one since \(M\) is self-dual.

It follows that any of its vectors is a linear
combination of products of powers \(a_j\) of
prime-to-\(p\) translates of \(\theta\). The composition
scalars in such a sum are roots of unity. A further
prime-to-\(p\) translate \(t\) of the vector modifies
the one-variable translates by coordinates with
prime-to-\(p\) denominators and by the same type of
scalars.

\proofstep{Transport the expansion and evaluate its coordinates by the product Mahler measure.}
Transport this expansion for 
\[
  G=U_t e^\natural
\]
 and
apply the product Mahler measure in the \(x_j\) frame.
Its norm is at most one because the \(x_j\) are test
directions for the actual normalized vector. Let
\(b_j\) be a different frame of lattice vectors that
is a basis after localization at \(p\). The transformation
to measure coordinates \(A\in\Z_p^{\dim M}\), with
\(\dim M\) the rank over \(O\), for this
frame is obtained by replacing the measure coordinates
in the \(x_i\) chart by
\[
                  A'_i=\sum_j A_j\lambda(\beta(x_i,b_j)).
\tag{M3}\label{eq:tag-m3}
\]
This is an integral homeomorphism since the Gram
diagonal entries and the split change of basis are
units.

If \(v\) is an algebraic linear combination
of lattice directions, then differentiating in
the direction \(\delta_0^{-1}v\) gives the linear
factor
\[
         L_v(A)=\Omega_0^{-1}\sum_j
               A_j\lambda(\beta(v,b_j)).
\tag{M4}\label{eq:tag-m4}
\]
To see the formula, write \(v=\sum_i v_i x_i\), compute
\(\Omega_0^{-1}\sum_i\lambda(v_i)A'_i\) by
\((1+X_i)d/dX_i\), and substitute \eqref{eq:tag-m3}.

For a prime-to-\(p\) integer \(N_0\) with
\(N_0t\in M\), consider the primary translates generated
by \(N_0\eta x/p^n\), where \(x\in M\), \(\eta\in O\)
and \(\eta_-\in p^{2n}\Z_p\). Their sums have the same
divisibility in the minus factor. Choose a common
\(\epsilon\in N_0O\) with components \((1,0)\) modulo
\(p^{2n}\). Reducing the plus component in the \(b_j\)
frame expresses each such translate \(y\), modulo
\(N_0 M\), as 
\[
  \epsilon\sum_j y_j b_j/p^n.
\]
The difference in the minus component is already
divisible by \(p^n\). The difference in the plus
component is divisible by \(p^n\) by the definition
of the digits \(y_j\). Both the new representative
and \(y\) have integral \(\beta\)-pairing with the
difference, since the original numerators have
minus component divisible by \(p^{2n}\).

For \(l\in N_0 M\) the operator \(W_l\) fixes \(G\),
because \(W_l\) commutes with \(U_t\) up to the
character \(\kappa(l,t)=1\).
In particular a change by \(l\in N_0M\) introduces the character factor
\[
  \kappa(l,t)=1,\qquad W_l U_tG=U_tW_lG\quad(l\in N_0M).
\]
 The composition law
therefore preserves \(W_yG\) under this change of
representative. The pairing with the measure
coordinates is preserved as well.

Let \(D\) be a common multiple of the prime-to-\(p\) denominators
of the tensor translates and of the \(a_i\), prime to \(p\).
The Chinese remainder theorem gives \(s\equiv0\bmod D\),
\(s\equiv1\bmod p^n\). This clears the coordinates for \eqref{eq:tag-m2} in
the \(x_i\) frame without changing the action, since
\[
 (s-1)y\in N_0M,\qquad \beta(y,y)\in\mathbb Z
       \quad\Longrightarrow\quad W_{sy}G=W_yG .
\]
If the denominators of the prime-to-\(p\) translates
are cleared by \(s_0\mid s\), each coordinate of
\(s y\) has the form \(s_0 v_i/p^n\) with
\(v_i\in O\) and \((v_i)_-\in p^{2n}\Z_p\).
That is the lift in \eqref{eq:tag-m2}.

In the product the
torsion character from this action is the image
under \(\lambda\) of
\(\kappa(\sum_j (A_j\bmod p^n)b_j,y)\).
Indeed on the split torsion the alternating pairing is
the pairing of the opposite coordinate of its first
argument with the plus coordinate of its second
argument. In the \(x_i\) frame this is the sum
of the one-variable pairings with weights \(a_i\).
Formula \eqref{eq:tag-m3} includes exactly these weights and the
conjugated first-argument coordinates. Reducing an
exponent modulo \(p^n\) leaves the character unchanged.

If \(\zeta^{p^n}=1\), this invariance follows from
\[
  m'=m+p^n h\quad\Longrightarrow\quad\zeta^{m'}=\zeta^m(\zeta^{p^n})^h=\zeta^m .
\]
Thus the derivatives of \(W_yG\) in a list of scaled
directions are the integrals of the products of
\eqref{eq:tag-m4} times this torsion character. The same measure,
of norm at most one, works for all such translates
and lists of directions. The transported derivatives
are algebraic by the one-variable argument, the
product rule and algebraic directions. In the tensor
expansion the scalar of a normalized vector is
algebraic by comparing one nonzero scaled jet.
If \(J\) is the nonzero scaled jet chosen for this comparison and \(G=cG_0\), then
\[
  c=\frac{J(G)}{J(G_0)}\qquad(J(G_0)\ne0).
\]
Such a jet is algebraic for the actual vector by
its bound one for all field transports, and for the
expansion by the same one-variable argument.
\label{cor:word-measure-6}

For the actual period reading, write \(b_j\) in complex
coordinates as \(\delta\operatorname{emb}(B_j)\), with
\(B_j\) rational Hom maps, and take the direction \(v\)
in \eqref{eq:tag-m4} equal to \(\delta_0\operatorname{emb}(V)\).
The factor \(\sqrt d/(2\operatorname{area}(\Gamma))\)
in \(\beta\) equals \((\abs{\delta}^2 N\mathfrak c)^{-1}\).
Consequently \eqref{eq:tag-m4} is the trace pairing with
\(\sum_j A_j B_j\), divided by the period
\[
       \Omega=\Omega_0\, N\mathfrak c\,\lambda(\delta/\delta_0),
                      \qquad \abs{\Omega}_p=1.
\tag{M5}\label{eq:tag-m5}
\]

In entries the trace pairing with a matrix unit is the
corresponding adjoint entry. The diagonal derivation
generators give the units
\(u_\tau=\lambda(c_{D,\tau})/\Omega\), whose unit
valuations are the chosen derivation normalization at the
prime assigned to the sequence. Hence the polynomial belonging to the jet
index \(m\) is exactly the monomial of the point
evaluation in Section~\ref{subsec:word-moments}.
Multilinearity applies to a repeated derivative of any
degree without a factorial denominator. Passing from the
analytic Taylor coefficient to the jet multiplies by the
multi-index factorial.
In coordinates the conversion is
\[
  \partial^m F(0)=\left(\prod_j m_j!\right)[X^m]F .
\]

\proofstep{Approximate the entire finite list simultaneously by residue balls.}
The characteristic functions of residue balls
in \(\Z_p^{\dim M}\) can now be used to approximate each
finite list of such integrands simultaneously. In this
Riemann sum, assign a ball the measure of its indicator
and sample it at a nonnegative integer point in the ball.
The norm bound makes every measure coefficient integral.
There are no losses as the precision or list grows.
The resulting sample is a rational point in the Hom
period lattice and has exactly the pairing monomial
and torsion character just computed.
\end{proof}

